\begin{abstract}
Cross-silo federated learning (FL) between hospitals is limited less by computation than by the uplink: every synchronisation round is a scheduled, audited transfer over an institutional WAN, a cellular edge slice or a satellite backhaul. We treat the learning job as a network tenant with an explicit communication quota and study, on a patient-deduplicated readmission benchmark with ten seeds per configuration and \num{TotalRuns} runs in total, how the number of rounds in which a site transmits governs three costs at once: accuracy, network cost, and the privacy budget under client-level differential privacy (DP). We introduce \tierfed{}, a scheduler that measures per-tier cross-client dispersion and solves, exactly and in under a millisecond, the integer programme that minimises a local-SGD drift bound subject to the quota, while keeping a per-site classifier head private. At a quota of $\rho=0.25$, \tierfed{} reaches \num{diabetesmlpTierfedAuroc}~AUROC against \num{diabetesmlpFedperAuroc} for full-rate FedPer, \ifflag{diabetesmlpTierfedNonInferiorAuroc}{non-inferior on AUROC}{not non-inferior on AUROC}\ifflag{diabetesmlpTierfedNonInferiorAuprc}{ and on AUPRC}{} at a 0.01 margin, with \num{diabetesmlpTierfedUploadSavingPct}\% less upload; composed with sparse ternary compression and error feedback it uploads \num{compTierfedStcUploadMb}~MB per site instead of \num{compFedperNoneUploadMb}~MB at \num{compTierfedStcAuroc}~AUROC. Under an identical client-level DP mechanism with an exact R\'enyi accountant, transmitting in fewer rounds is charged a smaller budget and injects noise fewer times: at noise multiplier $\sigma=1$, \tierfed{} reaches \num{dpTierfedSigmaOneAuroc}~AUROC at $\varepsilon=\num{dpTierfedSigmaOneEps}$ against \num{dpFedperSigmaOneAuroc} at $\varepsilon=\num{dpFedperSigmaOneEps}$ for FedPer. We also report what does not favour the method. A fixed uniform period and the AdaComm baseline match or exceed \tierfed{}'s accuracy at comparable upload without DP; the divergence-driven tier differentiation is rarely active on this data; no method separates from the others under update inversion or membership inference; and a claim of an earlier version of this work, that sparse multi-rate schedules are harmful, does not survive patient-level deduplication and is retracted. Every number in this article is generated from released, bit-reproducible per-seed logs.
\end{abstract}
